% =====================================================================
%  Chapter 15 — The annulus residual (S41)
% =====================================================================
\chapter{The Annulus: A Route Proved Impossible}
\label{ch:annulus}

\begin{record}
\textbf{What this chapter covers.} Session \Sn{41}. It was commissioned to
close the last technical item in H1's estimate layer --- an item \Sn{36} had
classified as \emph{open-mechanical}, ``bookkeeping\dots{} it involves no new
estimate''. It returned a proof that the proposed method \emph{cannot} close
it, for any choice of parameters, together with a second, independent proof
that no variant of the method can either.

\smallskip
The session is the programme's cleanest example of its own reporting rule.
It was asked for a closure and delivered an impossibility; the impossibility
is logged as the main finding, and the item it concerns was reclassified from
``mechanical'' to genuinely \OPEN{} --- a strictly worse position than the
programme believed it was in. It is also the session that established the
merge gate of Chapter~\ref{ch:house-rules}, which caught the inverted lemma
three sessions later.
\end{record}

\section{What was to be closed}
\label{sec:annulus-setting}

\texttt{prop:K-amended} (line~6215) bounds the \(K\)-coupling with an absolute
constant --- no logarithm, no correlation hypothesis --- modulo the annulus
residual \(\Aann\) of Definition~\ref{def:aann}: the cutoff-derivative terms
supported on the transition band \(\{M/8\le\abs{\omega}\le M/4\}\), carrying
\emph{incomplete} weights \(\psi\abs{\chi'}\) instead of the full \(\eta^2\).

\texttt{rem:level-iteration} (line~6249) proposed the standard De Giorgi
device. Replace the single level pair \((\tfrac18,\tfrac14)\) by a nested
sequence \(\theta_k=\tfrac18+2^{-k-3}\) with cutoffs \(\chi_k\) at levels
\((\theta_{k+1},\theta_k)\). Each level-\(k\) annulus term is dominated by the
full-weight quantities at level \(k+1\) with constant \(C4^k\), and the
iteration is closed by choosing Young parameters \(\delta_k=\delta_08^{-k}\).
The principle invoked was ``geometric beats geometric''.

Nobody had carried it out.

\section{The positive half: an exact inventory}
\label{sec:inventory}

Before attacking the iteration, \Sn{41} enumerated what actually has to be
absorbed. The weighted magnitude identity produces exactly eight families of
\(\chi\)-derivative terms, and \texttt{prop:s41-core} (line~8011) classifies
all eight.

\begin{table}[ht]
\centering
\scriptsize
\begin{tabular}{@{}l >{\raggedright\arraybackslash}p{0.13\textwidth} >{\raggedright\arraybackslash}p{0.46\textwidth} >{\raggedright\arraybackslash}p{0.17\textwidth}@{}}
\toprule
Family & Source & Form & Disposal \\
\midrule
G1 & \(R_3\)
  & \(-\tfrac{2\nu}{M}\iint m\,w\abs{\nabla m}^2\psi^2\chi\chi'\)
  & good sign, discarded \\
G2 & \(R_5\), sink \(-\nu m w\)
  & \(-\tfrac{\nu}{M}\iint m^3w^2\psi^2\chi\chi'\)
  & good sign, discarded \\
G3 & \(R_5\), \(\nu\Delta m\), \((\chi')^2\) piece
  & \(-\tfrac{\nu}{M^2}\iint m^2w\abs{\nabla m}^2\psi^2(\chi')^2\)
  & good sign, discarded \\
G4 & same, \(\nabla(m^2)\) piece
  & \(-\tfrac{2\nu}{M}\iint m\,w\abs{\nabla m}^2\psi^2\chi\chi'\)
  & good sign, discarded \\
\addlinespace
I1 & \(R_5\), through \(\alpha m\)
  & \(\tfrac1M\iint\alpha m^3w\,\psi^2\chi\chi'\),
    \(\le C\Lfac^2M^{3/2}\nu^{-1}\)
  & allowed inhomogeneous class \\
I2 & \(R_5\), through \(-m\dot M/M\)
  & same bound, via \(\abs{\dot M}/M\le\norm{\nabla u}_{L^\infty}\le CM\Lfac\)
  & allowed inhomogeneous class \\
\addlinespace
\textbf{S1} & \(\chi\chi''\) piece of G3's integration by parts
  & \(-\tfrac{\nu}{M^2}\iint m^2w\abs{\nabla m}^2\psi^2\chi\chi''\),
    \emph{indefinite sign}
  & \textbf{survives} \\
\textbf{S2} & cross family, from \(R_5\) and \(R_4\)
  & \(\tfrac{c\nu}{M}\iint m^2\psi^2\chi\chi'\,\nabla m\cdot\nabla w\),
    \(\abs{c}\le2\), \emph{indefinite sign}
  & \textbf{survives} \\
\bottomrule
\end{tabular}
\caption[The eight cutoff-derivative families]{The complete inventory of
\texttt{prop:s41-core} (line~8011). Six of the eight families are disposed of
outright. The irreducible core is exactly two, both of indefinite sign, and
the negative results of \S\ref{sec:noclose} are all about them.}
\label{tab:s41-families}
\end{table}

This inventory is worth its own emphasis. It is a genuine positive result of a
session whose headline is negative: three quarters of the residual is not
there at all, and everything the programme now has to argue about is
\(\abs{\mathrm{S1}}+\abs{\mathrm{S2}}\).

\section{The negative half: the iteration cannot close}
\label{sec:noclose}

\subsection{No contraction, for any parameter}

Executing the sketched iteration on the core gives a recursion
\(Y_k\le G_k+\Lambda_k Y_{k+1}\) with
\begin{equation}\label{eq:s41-lambda}
  \Lambda_k \;=\; 2^{8}\cdot4^{k} \;+\; 2^{8}\,\delta_k^{-1}\cdot4^{k}.
\end{equation}

\begin{statement}[\texttt{prop:s41-nocontraction}, line~8189]
\label{st:nocontraction}
\IMPOSSIBLE{} \(\Lambda_k\ge2^{8}\cdot4^{k}\ge2^{8}>1\) for \emph{every}
sequence \((\delta_k)\). The recursion never contracts.
\end{statement}

The reason Young's inequality cannot help is visible in
\eqref{eq:s41-lambda}: the first term does not involve \(\delta_k\) at all.
The S1 family is not a product of a dissipative factor with anything else, so
there is no place to insert a small parameter. Three consequences follow
immediately and each is recorded:

\begin{itemize}[leftmargin=1.6em]
\item the originally proposed \(\delta_k=\delta_08^{-k}\) makes matters
  \emph{strictly worse}: \(\Lambda_k\sim32^k\) rather than \(4^k\);
\item the summability \(\sum_k\delta_k\le\tfrac12\) that absorption
  independently requires forces \(\delta_k^{-1}\ge2\), so the two geometric
  factors land on the same side and reinforce;
\item ``geometric beats geometric'' has no realization here.
\end{itemize}

\subsection{The unrolled form is worse than trivial}

\begin{statement}[\texttt{prop:s41-vacuous}, line~8224]
\label{st:vacuous}
\IMPOSSIBLE{} The unrolled inequality is implied by \(Y_0\le Y_0\) and carries
no information, for every \(N\); letting \(N\to\infty\) does not help. The
remainder coefficient \(\prod_i\Lambda_i\) diverges while \(Y_N\ge Y_0\) by
monotonicity.
\end{statement}

\subsection{Why no cleverer choice of levels exists}

The two results above concern one scheme. The third rules out the family.

\begin{statement}[\texttt{prop:s41-coarea}, line~8303]
\label{st:coarea}
\IMPOSSIBLE{} As \(k\to\infty\) the nested cutoffs converge weakly-\(*\) to a
Dirac mass at the boundary level, so the residual converges to a literal
\emph{flux through the level set} \(\{\abs{\omega}=M/8\}\) --- which is exactly
the pointwise information H1 is trying to produce in the first place. More
generally, by the scale-free normalisation \(\int\abs{\chi'}\dee s=1\),
\emph{any} admissible level cutoff produces a residual that is a level density
of the underlying measure \(\mu_F\), bounded above only by
\((\theta_+-\theta_-)^{-1}\le16\) times a total mass, and bounded below not at
all.
\end{statement}

A pigeonhole check makes the last point concrete and is worth reproducing
because it explains why redistributing levels is futile: with \(N\)
equally-spaced levels, \(\abs{\chi'}\le16N\) runs against a best-annulus mass
\(\le1/N\), and the two cancel exactly. Uniform, geometric or optimally chosen
levels all land in the same place.

\subsection{The root cause}

\texttt{rem:s41-whichlemma} (line~8249) identifies why neither classical
iteration lemma is available, and this is the part of the session that
generalises.

\begin{itemize}[leftmargin=1.6em]
\item The \emph{linear} lemma --- Ladyzhenskaya--Ural'tseva, or
  Giaquinta--Giusti hole-filling --- needs a contraction factor below one,
  which Statement~\ref{st:nocontraction} proves impossible.
\item The \emph{nonlinear} lemma --- Stampacchia--De Giorgi --- needs the
  coarse level bounded by the fine one with a superlinear gain and small
  initial data. Three failures occur simultaneously: the recursion runs
  coarse-to-fine, the opposite of standard De Giorgi truncation; it is
  homogeneous of degree one, so there is no superlinear gain; and smallness of
  \(Y_0\) is the sought \emph{conclusion}, not an available hypothesis.
\end{itemize}

\begin{obsv}[Why De Giorgi truncation does not apply here]
\label{obs:degiorgi}
De Giorgi truncation gains because the truncated sets \emph{shrink}: superlevel
sets contract to measure zero, and the nonlinear gain feeds on that decay of
measure. In this problem the truncation is not in the solution but in an
\emph{auxiliary} variable, \(\abs{\omega}\), and it is confined to a fixed band
\(\{M/8\le\abs{\omega}\le M/4\}\). No set shrinks. No measure decays. There is
nothing for a nonlinear gain to feed on.

\smallskip
This is a statement about the geometry of the problem and not about the
skill of the execution, which is why \texttt{prop:s41-coarea} can rule out
every level selection at once rather than one at a time.
\end{obsv}

\begin{figure}[ht]
\centering
\begin{tikzpicture}[
  >={Latex[length=2mm]},
  font=\scriptsize
]
% --- left panel: classical De Giorgi ---
\begin{scope}
\node[font=\small\bfseries] at (2.1,3.3) {Classical De Giorgi};
\foreach \k/\yb/\yt/\op in {0/0.0/2.4/12, 1/0.45/2.05/26, 2/0.85/1.7/40,
                            3/1.15/1.45/58}{
  \fill[provedgreen, opacity=0.\op] (0,\yb) rectangle (4.2,\yt);
}
\draw[thick] (0,0) rectangle (4.2,2.4);
\node[align=center, text width=40mm] at (2.1,-0.75)
  {superlevel sets \emph{shrink} to measure zero;\\
   the nonlinear gain feeds on that decay};
\draw[->, thick, provedgreen] (4.5,1.2) -- (5.15,1.2)
  node[right, text=provedgreen] {closes};
\end{scope}
% --- right panel: the annulus band ---
\begin{scope}[xshift=7.6cm]
\node[font=\small\bfseries] at (2.1,3.3) {The annulus residual};
\fill[panelgrey] (0,0) rectangle (4.2,2.4);
\fill[impossiblered, opacity=0.18] (0,0.85) rectangle (4.2,1.55);
\foreach \k in {0,1,2,3}{
  \draw[impossiblered, thin] (0,0.85+0.09*\k) -- (4.2,0.85+0.09*\k);
  \draw[impossiblered, thin] (0,1.55-0.09*\k) -- (4.2,1.55-0.09*\k);
}
\draw[thick] (0,0) rectangle (4.2,2.4);
\node[right, font=\scriptsize] at (4.25,1.55) {\(\abs{\omega}=M/4\)};
\node[right, font=\scriptsize] at (4.25,0.85) {\(\abs{\omega}=M/8\)};
\node[align=center, text width=44mm] at (2.1,-0.85)
  {truncation is in the \emph{auxiliary} variable \(\abs{\omega}\),
   inside a \emph{fixed} band: nesting the levels moves nothing,
   and no set shrinks};
\draw[->, thick, impossiblered] (-0.75,1.2) -- (-0.1,1.2);
\node[left, text=impossiblered, align=right, text width=17mm] at (-0.8,1.2)
  {\(\Lambda_k\ge2^8\)\\ always};
\end{scope}
\end{tikzpicture}
\caption[Why the level iteration cannot work]{The structural difference
\texttt{rem:s41-whichlemma} (line~8249) identifies. On the left, the mechanism
that makes De Giorgi iteration work. On the right, the configuration
\(\Aann\) actually presents: the levels are nested inside a band of fixed
relative width, so refining them buys no measure decay, and by
\texttt{prop:s41-coarea} the limit of the nesting is a flux through a level
set --- the very object the argument was trying to produce.}
\label{fig:annulus}
\end{figure}

\section{\texorpdfstring{The shortfall: \texttt{rem:s41-shortfall}}{The shortfall}}
\label{sec:s41-shortfall}

What survives unconditionally is a crude bound.
\texttt{prop:s41-crude} (line~8395) gives, in the \(K\)-units of
\texttt{prop:crude-K},
\[
  M^{-2}\,\Aann^{\mathrm{core}}
  \;\le\; C\left(1+\delta^{-1}\right)\Lfac\,\norm{w}_{L^\infty}(\rstar)^3
    \;+\;\delta D_\flat .
\]

\begin{record}
\textbf{\texttt{rem:s41-shortfall}} (line~8428).
``\texttt{prop:s41-crude} reproduces, term for term, the bound of
\texttt{prop:crude-K}. The quantity that \texttt{prop:second-rung} is able to
absorb is \(\delta\nu(\rstar)^3\norm{w}_{L^\infty}\); what
\texttt{prop:s41-crude} delivers is \(C\Lfac\norm{w}_{L^\infty}(\rstar)^3\).
Absorption therefore requires \(C\Lfac\le\delta\nu\), which fails for fixed
\(\nu\) as \(\Lfac\to\infty\). So the annulus residual is short of closure by
exactly the same single logarithm that \texttt{prop:crude-K} identified for
\(K\) itself --- and, by \texttt{rem:s41-reading}, no level scheme recovers
it.''
\end{record}

This is the second of the Calder\'on--Zygmund logarithm's three appearances,
and \Sn{41} recognised it as a recurrence rather than a new difficulty --- but
related it to \Sn{34}, not to Defect~D8 of the audit, which is where it
originated. Chapter~\ref{ch:cz-log} follows the whole chain, and
Chapter~\ref{ch:inheritance} explains why the connection to the audit took
sixteen sessions to make.

\section{The replacement target}
\label{sec:band-target}

\texttt{rem:s41-endtarget} (line~8442) states exactly what would have to be
proved for \texttt{prop:K-amended}'s ``modulo \(\Aann\)'' caveat to be
removable, and stresses two points that a later session could easily get
wrong.

First, the absorbed term must carry \emph{the same} cutoff \(\eta\) as the
left-hand side. An annulus bound in terms of a fuller weight --- \(\psi^2\) on
the band, or \(\eta_{k+1}^2\) --- does not absorb, ``and this, not the value of
any constant, is what defeats the nested scheme''.

Second, by \texttt{lem:s41-admissible} (line~7830) it suffices to prove the
target for a \emph{single} admissible level pair, because all such pairs are
interchangeable downstream. The difficulty is therefore genuinely analytic and
not a matter of choosing levels cleverly --- which is the same conclusion the
co-area argument reaches from the other side.

The statement is then packaged as \texttt{target:band-absorption} (line~8468),
``the replacement for \texttt{rem:level-iteration}'': there exist levels
\(\theta_-<\theta_+\) in \([\tfrac1{16},\tfrac14]\), a cutoff \(\chi\) and a
constant \(C\) such that, with \(\Band=\{\theta_-M\le m\le\theta_+M\}\cap\operatorname{supp}\psi\),
\begin{align}
  \nu\iint_{\Band}\abs{\nabla m}^2w\,\psi^2
    &\le 2^{-11}\nu\iint\abs{\nabla m}^2w\,\eta^2
      + CM^2\nu\iint\abs{\nabla^2\ehat}^2\eta^2
      + C\Lfac^2M^{3/2}\nu^{-1},
      \label{eq:bandY}\\
  M^2\nu\iint_{\Band}\abs{\nabla^2\ehat}^2\psi^2
    &\le CM^2\nu\iint\abs{\nabla^2\ehat}^2\eta^2 + \cdots
      \label{eq:bandD}
\end{align}
Equation~\eqref{eq:bandY} is the \(Y\)-half and \eqref{eq:bandD} the
\(D\)-half. The distinction matters and is the subject of
Chapter~\ref{ch:non-implication}: they are not the same kind of statement, and
no hypothesis in the programme's family covers the second.

\section{The consequence for the ledger}
\label{sec:s41-ledger}

\begin{record}
\textbf{The reclassification.} \(\Aann\) moves from \emph{open-mechanical}
--- \Sn{36}'s belief that it was free bookkeeping --- to genuinely \OPEN{}.
H1 still rests on exactly two things: (a) \(\sstar\le\varepsilon_0\), and
(b) \texttt{target:band-absorption}. The \emph{count} is unchanged. What
changed is that (b) is now known to be real analytical content rather than a
formality.

\smallskip
The session log states the reading without softening: ``This is not progress
toward closing H1, and is reported as such; it is, however, a correction of a
false expectation, and the honest thing to log.''
\end{record}

Nothing else moved. H2, \(\sstar\)-decay, Type-II exclusion, the withdrawn
\S20--\S22 material and the Prize and Claim-A conclusions were all untouched,
and the merge diff was confined to three locations: the superseded block
inside \texttt{rem:level-iteration}, one row of the \Sn{36} status table, and
the new section.

\section{The session that built the gate}
\label{sec:s41-gate}

\Sn{41} is where the merge gate of Chapter~\ref{ch:house-rules} was
established, and the reasoning behind it was explicit: the programme's history
--- \Sn{31}'s arithmetically impossible induction, reversed Jensen inequality
and circularity; \Sn{35}'s retracted bonus term --- made a self-report an
unacceptable basis for merging into the shared document.

The subagent worked in an isolated \texttt{git} worktree and was instructed to
trace every symbol to its real definition rather than guess notation, to carry
out the convergent-series computation rather than assert it, to report a clean
failure if the argument did not close, and to stay scoped to one item. The
orchestrator then, before merging:

\begin{itemize}[leftmargin=1.6em]
\item checked brace balance programmatically (final depth zero, never
  negative) and the \(\backslash\)\texttt{begin}/\(\backslash\)\texttt{end}
  multiset;
\item confirmed every environment used was standard or declared, flagging two
  pre-existing undeclared environments as \emph{not} the session's own and
  verifying that they predated the diff;
\item confirmed via \texttt{git diff --stat} that exactly one file changed,
  \(864\) insertions and \(6\) deletions, all in three locations;
\item \textbf{re-derived the two load-bearing computations by hand} --- the
  Young step \(4XZ\le\delta X^2+4\delta^{-1}Z^2\), and the formula
  \eqref{eq:s41-lambda} for \(\Lambda_k\) --- rather than trusting them;
\item read the whole new section once as a sceptical referee, watching for the
  specific error classes previous audits had caught.
\end{itemize}

Only then was the file copied out of the worktree. Three sessions later the
same procedure caught an inverted volume factor that would otherwise have
entered the document as good news. That episode is
Chapter~\ref{ch:non-implication}.

\begin{obsv}[The value of executing a sketch]
\label{obs:execute-the-sketch}
\texttt{rem:level-iteration} had stood in the document for five sessions,
labelled ``no new estimate'', and everything built on top of it was
provisionally sound on that basis. The cost of finding out otherwise was one
session. The cost of \emph{not} finding out would have been every subsequent
ledger in the programme carrying a false entry in the column that says how
much work remains.

\smallskip
This is the strongest practical argument in the book for the discipline of
Chapter~\ref{ch:house-rules}: a claim marked ``mechanical'' is a claim, and it
is the kind of claim least likely ever to be checked, because nobody
volunteers to verify bookkeeping.
\end{obsv}
